\documentclass[11pt,a4paper]{article}
\usepackage[a4paper,margin=0.5in]{geometry}
\usepackage[utf8]{inputenc}
\usepackage{mdwlist}
\usepackage[default]{lato}
\usepackage[T1]{fontenc}
\usepackage{textcomp}
\usepackage{fontawesome}
\usepackage{enumitem}
\usepackage{newunicodechar}
\newunicodechar{（}{(}
\newunicodechar{）}{)}
\DeclareFontFamily{U}{fontawesomeOne}{}
\DeclareFontShape{U}{fontawesomeOne}{m}{n}
{<-> FontAwesome--fontawesomeone}{}
\DeclareRobustCommand\FAone{\fontencoding{U}\fontfamily{fontawesomeOne}\selectfont}
\pagestyle{empty}
\setlength{\tabcolsep}{0em}

% format two pieces of text, one left aligned and one right aligned
\newcommand{\headerrow}[2]
{\begin{tabular*}{\linewidth}{l@{\extracolsep{\fill}}r}
#1 &
#2 \\
\end{tabular*}}

\begin{document}

\begin{center}
	{\LARGE \textbf{FENG-AN (Alan) HSIEH}}\\
	Zhubei, Hsinchu, Taiwan
	\vspace{0.05cm}
	\\
	\raisebox{-0.2\height}{\Large \faPhoneSquare} \ \ (+886) 986161907
    \raisebox{-0.2\height}{\Large \faEnvelopeSquare} \ \ xiefengan013@gmail.com
    \raisebox{-0.2\height}{\Large \faLinkedinSquare} \ \ linkedin.com/in/feng-an-hsieh-a0a69a15a/
\end{center}

\subsection*{\large Summary}
\hrule
\vspace{0.3em}

\hyphenpenalty=1000
\begin{itemize}[leftmargin=1em,noitemsep]
	\item[]
    Engineer with 4 years of C/C++ development at the hardware/firmware boundary: modeling hardware IP from RTL specifications, enabling firmware bring-up and debug on virtual platforms with processor fast models, and writing Python tooling and automation for firmware teams. Experienced in HW/SW co-simulation and in isolating bugs across HW and FW. Built AI-agent workflows that cut model delivery from $\sim$2 weeks to 2--3 days.
\end{itemize}

% ======================== Experience Section ==========================
\subsection*{\large Experience}
\hrule
\vspace{0.3em}

\renewcommand\labelitemi{}
\renewcommand\labelitemii{$\bullet$}
\begin{itemize}[leftmargin=1em]
	\parskip=0.1em

	\item
	      \headerrow
	      {\textbf{SiliconMotion} -- \textbf{ESL Modeling Team}}
	      {\textbf{Zhubei, Hsinchu}}
	      \headerrow
	      {\emph{Chief Engineer}}
	      {\emph{July 2026 -- Present}}
	      \begin{itemize*}
	      	\item Own 6 C++/SystemC TLM IP models written from RTL design specifications and delivered to 5+ firmware engineers, who use them to develop and debug firmware before silicon; models are also packaged into virtual platforms for external customers.
	      	\item Modeled controller data-path and control-path logic (command arbitration, DMA, error handling, command decode) and a shared inter-engine messaging IP that handshakes with every engine.
	      	\item Built an AI agent harness and reusable agent skills (Claude Code, Cursor) for spec-to-model scaffolding, documentation sync, and automated regression.
	      	\item Mentor engineers on TLM-2.0 modeling standards and development workflow.
	      \end{itemize*}
	      \headerrow
	      {\emph{Senior Engineer}}
	      {\emph{November 2022 -- June 2026}}
	      \begin{itemize*}
	      	\item Built and maintained full-system virtual platforms on Synopsys VDK, integrating processor fast models (Arm Cortex-R82, Cortex-R8, Cortex-A55; Andes D23) with peripheral TLM models, so firmware teams could bring up and test software months ahead of silicon.
	      	\item Ran HW/SW co-simulation with DV patterns and verified models against RTL, finding and resolving CModel and firmware bugs at the HW/FW boundary.
	      	\item Built Python tools for platform-level trace analysis and debugging, shortening issue isolation for firmware teams.
	      	\item Automated RTL waveform (FSDB/VCD) parsing with Python and shell, speeding up test creation by 3x and saving $\sim$1 week of manual effort per cycle.
	      	\item Profiled and tuned processor fast models and platform integration to improve simulation speed and stability.
	      \end{itemize*}
\end{itemize}

% ============================ Skill Section ===============================
\subsection*{\large Skills}
\hrule
\vspace{0.3em}

\hyphenpenalty=1000
\begin{itemize}[leftmargin=1em,noitemsep]
	\item \textbf{Embedded \& Firmware-facing:}
	      C, C++, Assembly, Arm Cortex-R/A and Andes cores (fast models), Firmware Debug on Virtual Platforms, HW/SW Co-simulation, Timer/Peripheral IP Modeling
	% TODO (add only if true): I2C / SPI / UART, CMIS, SerDes / Ethernet / PCIe / USB exposure
	\item \textbf{Scripting \& Tooling:}
	      Python (test automation, trace analysis, waveform parsing), Shell, Git
	\item \textbf{Modeling \& Verification:}
	      SystemC, TLM-2.0, Synopsys VDK, Verilog, RTL-vs-model Verification, FSDB/VCD
	\item \textbf{Protocols \& Architecture:}
	      AMBA spec \& bus integration (AXI, AHB, APB, CHI), DMA
    \item \textbf{Automation \& AI:}
          Claude Code, Cursor, Agent Skills, Regression Flows
    \item \textbf{Languages:}
          Mandarin (Native), English (TOEIC 780)
\end{itemize}

% ======================== Education Section ==========================
\subsection*{\large Education}
\hrule
\vspace{0.3em}

\begin{itemize}[leftmargin=1em]
	\parskip=0.1em

	\item
	      \headerrow
	      {\textbf{National Chung Cheng University}}
	      {\textbf{Chiayi, Taiwan}}
	      \headerrow
	      {\emph{M.E. in Electrical Engineering}}
	      {\emph{August 2020 -- July 2022}}
          \headerrow
          {\emph{Thesis: An End-To-End System For Road Agent Behavior Classification (Deep Learning, Computer Vision)}}
          {\emph{}}

	      \headerrow
          {\textbf{National University of Kaohsiung}}
	      {\textbf{Kaohsiung, Taiwan}}
          \headerrow
	      {\emph{B.E. in Computer Science and Information Engineering}}
	      {\emph{September 2016 -- June 2020}}
\end{itemize}

\end{document}
